\documentclass{beamer}
\usepackage[british,spanish]{babel}
\usepackage[utf8]{inputenc}

\usepackage{hyperref}
%\hypersetup{colorlinks=false,linkbordercolor=red,linkcolor=green,pdfborderstyle={/S/U/W 1}}

\usepackage{multirow}

\usepackage{textcomp}

\usepackage{listings}
\lstloadlanguages{Ruby}
\usepackage{cancel}

\usepackage{adjustbox}
\usepackage{lstcustom}

\usepackage{amsmath}

\usepackage{color}
\definecolor{light-gray}{gray}{0.80}
\definecolor{lstbackgroundshellcolor}{named}{light-gray}

\usepackage{tikz}
\newcommand*\circled[1]{\tikz[baseline=(char.base)]{
            \node[shape=circle,draw,inner sep=2pt] (char) {#1};}}

\usepackage[normalem]{ulem}

%\usepackage[acronym,xindy,toc]{glossaries}

\usepackage[acronym,xindy,toc]{glossaries}
\makeglossaries
%\usepackage[xindy]{imakeidx}
%\makeindex


\newcommand{\comment}[2]{#2}

\newcommand{\commandinline}[1]{\lstinline[basicstyle=\small\lstfontfamily]{#1}}
\newcommand{\outputcommand}[1]{\color{darkgreen}{#1}}

\graphicspath{ {./images/} }

\title{Ruby Basics}
%\subtitle[short subtitle]{long subtitle}
\author[C. Cuenca, F. Quintana]{Carmelo Cuenca-Hernández and Francisca Quintana-Domínguez}
%\institute{Escuela Universitaria de Informática}
%\date[04/2013]{Abril - 2013}
\date{}
\titlegraphic{\includegraphics[width=0.5 \textwidth]{./images/logo_ulpgc_version_horizontal_rgb.eps}}



\pgfdeclareimage[width=2.0\baselineskip]{ulpgc-logo}{./images/logosimbolo_secundario_version_vertical}
\setbeamertemplate{footline}{\raisebox{-2ex}{\pgfuseimage{ulpgc-logo}}
  \usebeamerfont{date in head/foot}\insertshortdate{}\hfill
  \usebeamertemplate{navigation symbols}\hfill
  \insertframenumber{}/\inserttotalframenumber}
\setbeamertemplate{sidebar right}{}


\usetheme{Antibes}
%\usetheme{Berlin}

%\usetheme{Warsaw}
%\usecolortheme{albatross}

\selectlanguage{british}



\begin{document}


\begin{frame}
	\titlepage
\end{frame}


\section*{Outline}
\begin{frame}[fragile, allowframebreaks]
  \frametitle{Outline}
  %\tableofcontents%[part=1,pausesections]
  %\tableofcontents[currentsection,currentsubsection, sectionstyle=show] 
  \tableofcontents[currentsection,sectionstyle=show,hideothersubsections]
\end{frame}

%%%%%%%%%%%%%%%%%%%%%%%%%%%%%%%%%%%%%%%%%%%%%%%%%%%%%%%%%%%%%%%%%%%%%%%%%%%%%%
%\newacronym{<label>}{<abbrv>}{<full>}
%\glsreset{<label>}
%\glsresetall
%\acrlong{<label>}
%\acrfull{<label>}
%\acrshort{<label>}
%\input{../glossary}

\newacronym{aws}{AWS}{Amazon Web Services}
\newacronym{css}{CSS}{cascading style sheets}
\newacronym{ebs}{EBS}{Elastic Block Storage}
\newacronym{ec2}{EC2}{Amazon Elastic Compute Cloud}
\newacronym{elb}{ELB}{Elastic Load Balancing}
\newacronym{ror}{RoR}{{\href{http://rubyonrails.org/}{Ruby on Rails}}}
\newacronym{rds}{RDS}{Relational Database Service}
\newacronym{rvm}{RVM}{{\href{https://rvm.io/}{Ruby Version Manager}}}
\newacronym{s3}{S3}{Simple Storage Service}
\newacronym{sqs}{SQS}{Amazon Simple Queue Service}
%%%%%%%%%%%%%%%%%%%%%%%%%%%%%%%%%%%%%%%%%%%%%%%%%%%%%%%%%%%%%%%%%%%%%%%%%%%%%%
\section{Ruby}
\begin{frame}[fragile]
\frametitle{Ruby}
\begin{columns}
 \column{0.20 \textwidth}{\includegraphics[width= 1.0 \textwidth]{ruby.eps}}
 \column{0.80 \textwidth}
 \begin{itemize}
  \item The inventor of Ruby, Yukihiro Matsumoto (1995) states in his book ``The Ruby Programming Language'' that:
  \begin{verbatim}
     It is easy to programs in Ruby,
     but Ruby is not a simple language.
  \end{verbatim} 
  \item Ruby is a dynamic, open source programming language with a focus on simplicity and productivity. It has an elegant syntax that is natural to read and easy to write

  \item \url{http://www.ruby-lang.org}
 \end{itemize}

\end{columns}

\end{frame}
%%%%%%%%%%%%%%%%%%%%%%%%%%%%%%%%%%%%%%%%%%%%%%%%%%%%%%%%%%%%%%%%%%%%%%%%%%%%%%
\section{Setting Up Ruby Environment}
\begin{frame}[fragile]
\frametitle{Setting Up Ruby Environment}
As user \texttt{invitado}, carry out the steps for installing Ruby 2.0.0

\lstset{language=shell}
\begin{lstlisting}[numbers=none, escapechar=!]
$ \curl -L https://get.rvm.io | bash -s stable --auto-dotfiles
!\outputcommand{A lot of stuff}!
$ source /home/invitado/.rvm/scripts/rvm
$ rvm install 2.0.0
$ rvm docs generate
\end{lstlisting}

\end{frame}
%%%%%%%%%%%%%%%%%%%%%%%%%%%%%%%%%%%%%%%%%%%%%%%%%%%%%%%%%%%%%%%%%%%%%%%%%%%%%%
%%%%%%%%%%%%%%%%%%%%%%%%%%%%%%%%%%%%%%%%%%%%%%%%%%%%%%%%%%%%%%%%%%%%%%%%%%%%%%
\section{Hello World!}
\begin{frame}[fragile]
\frametitle{Hello World!}
\begin{columns}
\column{0.5 \textwidth}
\lstset{language=Ruby, style=eclipse}
\begin{lstlisting}[escapechar=!]
# The famous Hello World
# Program is trivial in
# Ruby. Superfluous:
#
# * A "main" method
# * Newline
# * Semicolons
#
# Here is the Code:

puts "Hello World!"
\end{lstlisting}
\column{0.5 \textwidth}
\lstset{language=C, style=eclipse}
\begin{lstlisting}[escapechar=!]
#include<stdio.h>

int main(int arg_count,
         char ** arg_values)
{
 printf("Hello World\n");
 return 0;
}
\end{lstlisting}
\end{columns}
\begin{itemize}
 \item Use your favorite editor to open a new file with the filename \texttt{hello-world}.rb and insert the above lines into it. You can execute this Ruby program in the command lines as follows:

\lstset{language=shell}
\begin{lstlisting}[ numbers=none, escapechar=!]
$ ruby hello-world.rb
\end{lstlisting}
\end{itemize}



\end{frame}
%%%%%%%%%%%%%%%%%%%%%%%%%%%%%%%%%%%%%%%%%%%%%%%%%%%%%%%%%%%%%%%%%%%%%%%%%%%%%%
%%%%%%%%%%%%%%%%%%%%%%%%%%%%%%%%%%%%%%%%%%%%%%%%%%%%%%%%%%%%%%%%%%%%%%%%%%%%%%
\section{\texttt{irb}}
\begin{frame}[fragile]
\frametitle{\texttt{irb}}
\begin{itemize}
 \item \texttt{irb} stands for ``InteractiveRuby'' and is a kind of sandbox where you can play around with Ruby at your leisure, \texttt{irb} is launched by entering \texttt{irb} on the shell and ends if you enter \texttt{exit}
\lstset{language=shell}
\begin{lstlisting}[numbers=none,escapechar=&]
$ irb --simple-prompt
>> puts "Hello World!"
Hello World!
=> nil
>> exit
\end{lstlisting}
\end{itemize}
\end{frame}
%%%%%%%%%%%%%%%%%%%%%%%%%%%%%%%%%%%%%%%%%%%%%%%%%%%%%%%%%%%%%%%%%%%%%%%%%%%%%%
%%%%%%%%%%%%%%%%%%%%%%%%%%%%%%%%%%%%%%%%%%%%%%%%%%%%%%%%%%%%%%%%%%%%%%%%%%%%%%
\section{Ruby is Object Oriented}
\begin{frame}[fragile]
\frametitle{Ruby is Object Oriented}
\begin{itemize}
 \item Ruby only knows objects
 \item Everything is an object
 \item Every object is an instance of a class
 \item An object is encapsulated and can only be reached from the outside via the methods of the corresponding object.
 \item I cannot change any property of an object directly from the outside. The corresponding object has to offer a method with which I can do so
\end{itemize}

\end{frame}
%%%%%%%%%%%%%%%%%%%%%%%%%%%%%%%%%%%%%%%%%%%%%%%%%%%%%%%%%%%%%%%%%%%%%%%%%%%%%%
%%%%%%%%%%%%%%%%%%%%%%%%%%%%%%%%%%%%%%%%%%%%%%%%%%%%%%%%%%%%%%%%%%%%%%%%%%%%%%
\subsection{Methods}
\begin{frame}[fragile, allowframebreaks]
\frametitle{Methods}
\begin{block}{Important}
\alert{Names of methods are always written in snake case}
\end{block}

In other programming languages, the terms you would use for Ruby methods would be: functions, procedures, subroutines and of course methods

\begin{itemize}

\item We want to output the string \texttt{``Hello World!''} three times in a row
\lstset{language=Ruby, style=eclipse}
\begin{lstlisting}[escapechar=&]
# three_times.rb
def three_times
  puts 'Hello World!'
  puts 'Hello World!'
  puts 'Hello World!'
end

three_times
\end{lstlisting}


\lstset{language=shell}
\begin{lstlisting}[numbers=none, escapechar=&]
$ ruby three_times.rb
&\outputcommand{Hello World!\\
Hello World!\\
Hello World!}&
\end{lstlisting}

\item You can define parameters and use then within the method

\lstset{language=Ruby, style=eclipse}
\begin{lstlisting}[escapechar=&]
# three_times.rb
def three_times value
  puts value
  puts value
  puts value
end

three_times 'Hello World!'
\end{lstlisting}


\lstset{language=shell}
\begin{lstlisting}[numbers=none, escapechar=&]
$ ruby three_times.rb
&\outputcommand{Hello World!\\
Hello World!\\
Hello World!}&
\end{lstlisting}

\item You can give the variable \texttt{value} a default value and then you can also call the method without parameter
\lstset{language=Ruby, style=eclipse}
\begin{lstlisting}[escapechar=&]
# three_times.rb
def three_times value='Hello World!'
  puts value
  puts value
  puts value
end

three_times
\end{lstlisting}


\lstset{language=shell}
\begin{lstlisting}[numbers=none, escapechar=&]
$ ruby three_times.rb
&\outputcommand{Hello World!\\
Hello World!\\
Hello World!}&
\end{lstlisting}

\end{itemize}

\end{frame}
%%%%%%%%%%%%%%%%%%%%%%%%%%%%%%%%%%%%%%%%%%%%%%%%%%%%%%%%%%%%%%%%%%%%%%%%%%%%%%
\begin{frame}[fragile]
\frametitle{return}
\begin{itemize}
\item Ruby returns the value of the last expression as a default
\begin{columns}
 \column{0.5 \textwidth}
  \lstset{language=Ruby, style=eclipse}
  \begin{lstlisting}[escapechar=&]
def area_of_a_circle radius
  pi = 3.14
  return area = pi * radius * radius
end

puts area_of_a_circle 10
  \end{lstlisting}

 \column{0.5 \textwidth}
\lstset{language=Ruby, style=eclipse}
  \begin{lstlisting}[escapechar=&]
def area_of_a_circle radius
  pi = 3.14
  area = pi * radius * radius
end

puts area_of_a_circle 10
  \end{lstlisting}
\end{columns}
\end{itemize}

\lstset{language=shell}
\begin{lstlisting}[numbers=none, escapechar=&]
$ ruby area_of_a_circle
&\outputcommand{314.0}&
\end{lstlisting}

\end{frame}
%%%%%%%%%%%%%%%%%%%%%%%%%%%%%%%%%%%%%%%%%%%%%%%%%%%%%%%%%%%%%%%%%%%%%%%%%%%%%%
\subsection{Classes}
\begin{frame}[fragile, allowframebreaks]
\frametitle{Classes}
\begin{block}{Important}
\alert{The name of a class always starts with an upper case letter (Camel Case)}
\end{block}
\begin{itemize}
\item A class is a collection of methods

\lstset{language=Ruby, style=eclipse}
\begin{lstlisting}[escapechar=&]
# camel_case.rb
class CamelCase
  def three_times
    puts 'Hello World!'
    puts 'Hello World!'
    puts 'Hello World!'
  end
end

snake_case = CamelCase.new
snake_case.three_times
\end{lstlisting}

\lstset{language=shell}
\begin{lstlisting}[numbers=none, escapechar=&]
$ ruby three_times.rb
&\outputcommand{Hello World!\\
Hello World!\\
Hello World!}&
\end{lstlisting}

\end{itemize}
\end{frame}
%%%%%%%%%%%%%%%%%%%%%%%%%%%%%%%%%%%%%%%%%%%%%%%%%%%%%%%%%%%%%%%%%%%%%%%%%%%%%%
\begin{frame}[fragile, allowframebreaks]
\frametitle{Private methods}
\begin{itemize}
\item Quite often it make sense to only call a method within its own class or own instance.
Such methods are referred to as private methods

\lstset{language=Ruby, style=eclipse}
\begin{lstlisting}[escapechar=&]
# example.rb
class Example
  def a
    puts 'a'
  end

  private
  def b
    puts 'b'
  end
end

test = Example.new
test.a
test.b
\end{lstlisting}

\lstset{language=shell}
\begin{lstlisting}[numbers=none, escapechar=!]
$ ruby example.rb
!\outputcommand{a\\
example.rb:15:in `<main>': private method `b' called for \#<Example:0xa199b3c> (NoMethodError)}!
\end{lstlisting}
\end{itemize}
\end{frame}
%%%%%%%%%%%%%%%%%%%%%%%%%%%%%%%%%%%%%%%%%%%%%%%%%%%%%%%%%%%%%%%%%%%%%%%%%%%%%%
\begin{frame}[fragile, allowframebreaks]
\frametitle{Method \texttt{initialize()}}
\begin{itemize}
\item If a new instance is created, the method that is processed first and automatically is the method \texttt{initialize}.
The method is automatically a private method

\lstset{language=Ruby, style=eclipse}
\begin{lstlisting}[escapechar=&]
# room.rb
class Room
  def initialize
    puts 'abc'
  end
end

kitchen = Room.new
kitchen.initialize
\end{lstlisting}

\lstset{language=shell}
\begin{lstlisting}[numbers=none, escapechar=!]
$ ruby room.rb
abc
!\outputcommand{room.rb:9:in `<main>': private method `initialize' called for \#<Room:0x8529ce0> (NoMethodError)}!
\end{lstlisting}
\end{itemize}
\end{frame}
%%%%%%%%%%%%%%%%%%%%%%%%%%%%%%%%%%%%%%%%%%%%%%%%%%%%%%%%%%%%%%%%%%%%%%%%%%%%%%
\begin{frame}[fragile, allowframebreaks]
\frametitle{Inheritance}
\begin{itemize}
\item A class can inherit from another class. When defining the class, the parent class must be added with a \texttt{\textless} (smaller than) sign: 
\lstset{language=Ruby, style=eclipse}
\begin{lstlisting}[escapechar=&]
# a_b.rb
class A
  def a
    puts 'a'
  end
end

class B < A
  def b
    puts 'b'
  end
end

a = A.new
a.a
b = B.new
b.a
b.b
a.b
\end{lstlisting}

\lstset{language=shell}
\begin{lstlisting}[numbers=none, escapechar=!]
$ ruby a_b.rb
!\outputcommand{a\\
a\\
b\\
a\_b.rb:21:in `<main>': undefined method `b' for \#<A:0x9b258f0> (NoMethodError)\\}!
\end{lstlisting}
\end{itemize}
\end{frame}
%%%%%%%%%%%%%%%%%%%%%%%%%%%%%%%%%%%%%%%%%%%%%%%%%%%%%%%%%%%%%%%%%%%%%%%%%%%%%%
\begin{frame}[fragile, allowframebreaks]
\frametitle{Class Methods and Instance Methods}
\begin{itemize}
\item A class method can only be called in connection with the class
\item An instance method is a method than only works with an instance (you cannot apply the method \texttt{new} to an instance)


\lstset{language=Ruby, style=eclipse}
\begin{lstlisting}[escapechar=&]
# knowledge.rb
class Knowledge
  def pi
    puts 3.14
  end
end
knowledge = Knowledge.new
knowledge.pi
Knowledge.pi
\end{lstlisting}

\lstset{language=shell}
\begin{lstlisting}[numbers=none, escapechar=!]
$ ruby knowledge.rb
!\outputcommand{3.14\\
knowledge.rb:9:in `<main>': undefined method `pi' for Knowledge:Class (NoMethodError)
}!
\end{lstlisting}


\lstset{language=Ruby, style=eclipse}
\begin{lstlisting}[escapechar=&]
# knowledge.rb
class Knowledge
  def self.pi
    puts 3.14
  end
end
Knowledge.pi
knowledge = Knowledge.new
knowledge.pi
\end{lstlisting}

\lstset{language=shell}
\begin{lstlisting}[numbers=none, escapechar=!]
$ ruby knowledge.rb
!\outputcommand{3.14\\
knowledge.rb:9:in `<main>': undefined method `pi' for \#<Knowledge:0x9451c04> (NoMethodError)
}!
\end{lstlisting}

\item There are different notations for defining class methods

\lstset{language=Ruby, style=eclipse}
\begin{lstlisting}[escapechar=&]
# knowledge.rb
class Knowledge
  class << self
    def pi
      puts 3.14
    end
  end
end
Knowledge.pi
\end{lstlisting}
\end{itemize}
\end{frame}
%%%%%%%%%%%%%%%%%%%%%%%%%%%%%%%%%%%%%%%%%%%%%%%%%%%%%%%%%%%%%%%%%%%%%%%%%%%%%%
\section{Variables}
\begin{frame}[fragile]
\frametitle{Variables}
\begin{itemize}
 \item You already know that everything in Ruby is an object. So a variable must also be an object
\end{itemize}

\end{frame}
%%%%%%%%%%%%%%%%%%%%%%%%%%%%%%%%%%%%%%%%%%%%%%%%%%%%%%%%%%%%%%%%%%%%%%%%%%%%%%
\subsection{Naming Conventions}
\begin{frame}[fragile]
\frametitle{Naming Conventions}
\begin{block}{Important!}
\begin{itemize}
  \item In general variables are written in lower case. Constants start with an upper case letter
  \item A constant can also be overwritten with a new value in Ruby (Do not rely on the constancy of a constant!)
\end{itemize}

\end{block}

\end{frame}
%%%%%%%%%%%%%%%%%%%%%%%%%%%%%%%%%%%%%%%%%%%%%%%%%%%%%%%%%%%%%%%%%%%%%%%%%%%%%%
%%%%%%%%%%%%%%%%%%%%%%%%%%%%%%%%%%%%%%%%%%%%%%%%%%%%%%%%%%%%%%%%%%%%%%%%%%%%%%
\subsection{Strings}
\begin{frame}[fragile, allowframebreaks]
\frametitle{Strings}
\begin{itemize}
 \item Single Quotation Marks
\lstset{language=Ruby, style=eclipse}
\begin{lstlisting}[escapechar=&]
# string.rb
a = 'I am a string'
puts a
b = String.new
puts b
c = String.new 'I am also a string'
puts c
\end{lstlisting}
\lstset{language=shell}
\begin{lstlisting}[numbers=none, escapechar=!]
$ ruby string.rb
!\outputcommand{I am a string\\
\hspace{1cm}\\
I am also a string\\
}!
\end{lstlisting}
 \item Double Quotation Marks
\lstset{language=Ruby, style=eclipse}
\begin{lstlisting}[escapechar=&]
# string.rb
a = 'blue'
puts "Color #{a}"
\end{lstlisting}
\lstset{language=shell}
\begin{lstlisting}[numbers=none, escapechar=!]
$ ruby string.rb
!\outputcommand{
Color blue\\
}!
\end{lstlisting}

\end{itemize}
\end{frame}
%%%%%%%%%%%%%%%%%%%%%%%%%%%%%%%%%%%%%%%%%%%%%%%%%%%%%%%%%%%%%%%%%%%%%%%%%%%%%%
%%%%%%%%%%%%%%%%%%%%%%%%%%%%%%%%%%%%%%%%%%%%%%%%%%%%%%%%%%%%%%%%%%%%%%%%%%%%%%
\subsection{Integers}
\begin{frame}[fragile]
\frametitle{Integers}
\begin{itemize}
 \item A \texttt{Fixnum} is an \texttt{Integer} that can be saved in a ``word''. If a \texttt{Fixnum} gets bigger, it automatically becomes a \texttt{Bignum}
\lstset{language=Ruby, style=eclipse}
\begin{lstlisting}[escapechar=&]
# integer.rb
puts 33.class
a = 2014
puts a.class
a = a*9999
puts a.class
a = a*9999
puts a.class
\end{lstlisting}
\end{itemize}
\lstset{language=shell}
\begin{lstlisting}[numbers=none, escapechar=!]
$  ruby integer.rb 
!\outputcommand{Fixnum\\
Fixnum\\
Fixnum\\
Bignum
}!
\end{lstlisting}
\end{frame}
%%%%%%%%%%%%%%%%%%%%%%%%%%%%%%%%%%%%%%%%%%%%%%%%%%%%%%%%%%%%%%%%%%%%%%%%%%%%%%
%%%%%%%%%%%%%%%%%%%%%%%%%%%%%%%%%%%%%%%%%%%%%%%%%%%%%%%%%%%%%%%%%%%%%%%%%%%%%%
\subsection{Floats}
\begin{frame}[fragile]
\frametitle{Floats}
\begin{itemize}
 \item \texttt{Float} is a class for real numbers. The decimal separator is a point

\lstset{language=Ruby, style=eclipse}
\begin{lstlisting}[escapechar=&]
# float.rb
puts 3.14.class
pi = 3.14
puts pi
puts pi.class
\end{lstlisting}
\lstset{language=shell}
\begin{lstlisting}[numbers=none, escapechar=!]
$ ruby float.rb 
!\outputcommand{
Float\\
3.14\\
Float
}!
\end{lstlisting}
\end{itemize}

\end{frame}
%%%%%%%%%%%%%%%%%%%%%%%%%%%%%%%%%%%%%%%%%%%%%%%%%%%%%%%%%%%%%%%%%%%%%%%%%%%%%%
%%%%%%%%%%%%%%%%%%%%%%%%%%%%%%%%%%%%%%%%%%%%%%%%%%%%%%%%%%%%%%%%%%%%%%%%%%%%%
\subsection{Boolean Values and \texttt{nil}}
\begin{frame}[fragile]
\frametitle{Boolean Values and \texttt{nil}}
\begin{itemize}
 \item There are separate classes for boolean values (\texttt{true} and \texttt{false}) and for \texttt{nil} (no value) 
\lstset{language=Ruby, style=eclipse}
\begin{lstlisting}[escapechar=&]
# boolean.rb
puts true.class
puts false.class
puts nil.class
\end{lstlisting}
\lstset{language=shell}
\begin{lstlisting}[numbers=none, escapechar=!]
$ ruby boolean.rb 
!\outputcommand{
TrueClass\\
FalseClass\\
NilClass
}!
\end{lstlisting}
\item Number \texttt{0} is \texttt{true}
\item \texttt{nil} is false
\end{itemize}
\end{frame}
%%%%%%%%%%%%%%%%%%%%%%%%%%%%%%%%%%%%%%%%%%%%%%%%%%%%%%%%%%%%%%%%%%%%%%%%%%%%%%
%%%%%%%%%%%%%%%%%%%%%%%%%%%%%%%%%%%%%%%%%%%%%%%%%%%%%%%%%%%%%%%%%%%%%%%%%%%%%
\subsection{Scope of Variables}
\begin{frame}[fragile, allowframebreaks]
\frametitle{Scope of Variables}
\begin{itemize}
 \item Local variables start with a lower case or an underscore (\texttt{\_})
 \item A global variable starts with a \texttt{\$}-sign and is accessible in the entire program. They are used very rarely
 \item Instance variables or ``Attributes'' start with \texttt{@}) and they are only apply within a class
 \item Class variables are available throughout a class or module body and they start with  \texttt{@@}-signs
\end{itemize}
\lstset{language=Ruby, style=eclipse}
\begin{lstlisting}[escapechar=&]
# scope.rb

class Song
  @@count = 0

  def initialize title
    @title = title
    @@count += 1
  end

  def dump
    puts "title: #{@title}, count: #{@@count}"
  end
end

song1 = Song.new "No cambies"
song1.dump

song2 = Song.new "Mi carro"
song2.dump
\end{lstlisting}
\lstset{language=shell}
\begin{lstlisting}[numbers=none, escapechar=!]
$ ruby scope.rb 
!\outputcommand{
title: No cambies, count: 1\\
title: Mi carro, count: 2\\
}!
\end{lstlisting}
\end{frame}
%%%%%%%%%%%%%%%%%%%%%%%%%%%%%%%%%%%%%%%%%%%%%%%%%%%%%%%%%%%%%%%%%%%%%%%%%%%%%
\subsection{Getters and Setters}
\begin{frame}[fragile, allowframebreaks]
\frametitle{Getters and Setters}
\begin{itemize}
 \item As instance variable only exist within the relevant instance, you always need to write a ``getter'' method for exporting such a variable
\lstset{language=Ruby, style=eclipse}
\begin{lstlisting}[escapechar=&]
# getters_and_setters.rb

class Room
  def initialize
    @doors = 1
    @windows = 1
  end

  def doors
    @doors
  end

  def windows
    @windows
  end
end

room = Room.new
puts "doors: #{room.doors}, windows: #{room.windows}"
\end{lstlisting}
\lstset{language=shell}
\begin{lstlisting}[numbers=none, escapechar=!]
$ ruby getters_and_setters.rb
!\outputcommand{
doors: 1, windows: 1
}!
\end{lstlisting}
\item As this scenario is so common, there is already a ready-made getter method for it with the name \texttt{attr\_reader}
\lstset{language=Ruby, style=eclipse}
\begin{lstlisting}[escapechar=&]
# getters_and_setters.rb

class Room
  def initialize
    @doors = 1
    @windows = 1
  end

  attr_reader :doors, :windows
end

room = Room.new
puts "doors: #{room.doors}, windows: #{room.windows}"
\end{lstlisting}
\lstset{language=shell}
\begin{lstlisting}[numbers=none, escapechar=!]
$ ruby getters_and_setters.rb
!\outputcommand{
doors: 1, windows: 1
}!
\end{lstlisting}

\item A ``setter'' method changes an attribute from outside
\lstset{language=Ruby, style=eclipse}
\begin{lstlisting}[escapechar=&]
# getters_and_setters.rb

class Room
  def initialize
    @doors = 1
    @windows = 1
  end

  attr_reader :doors, :windows

  def doors= value
    @doors = value
  end

  def windows= value
    @windows = value
  end
end

room = Room.new
room.doors = 2
room.windows = 3
puts "doors: #{room.doors}, windows: #{room.windows}"
\end{lstlisting}

\lstset{language=shell}
\begin{lstlisting}[numbers=none, escapechar=!]
$ ruby getters_and_setters.rb
!\outputcommand{
doors: 2, windows: 3
}!
\end{lstlisting}

\item As this scenario is so common, there is already a ready-made setter method for it with the name \texttt{attr\_writer}
\lstset{language=Ruby, style=eclipse}
\begin{lstlisting}[escapechar=&]
# getters_and_setters.rb

class Room
  def initialize
    @doors = 1
    @windows = 1
  end

  attr_reader :doors, :windows
  attr_writer :doors, :windows
end

room = Room.new
room.doors = 2
room.windows = 3
puts "doors: #{room.doors}, windows: #{room.windows}"
\end{lstlisting}

\lstset{language=shell}
\begin{lstlisting}[numbers=none, escapechar=!]
$ ruby getters_and_setters.rb
!\outputcommand{
doors: 2, windows: 3
}!
\end{lstlisting}
\item There is even a method \texttt{attr\_accessor} that combines getter and setters
\lstset{language=Ruby, style=eclipse}
\begin{lstlisting}[escapechar=&]
# getters_and_setters.rb

class Room
  def initialize
    @doors = 1
    @windows = 1
  end

  attr_accessor :doors, :windows
end

room = Room.new
room.doors = 2
room.windows = 3
puts "doors: #{room.doors}, windows: #{room.windows}"
\end{lstlisting}

\lstset{language=shell}
\begin{lstlisting}[numbers=none, escapechar=!]
$ ruby getters_and_setters.rb
!\outputcommand{
doors: 2, windows: 3
}!
\end{lstlisting}


\end{itemize}

\end{frame}
%%%%%%%%%%%%%%%%%%%%%%%%%%%%%%%%%%%%%%%%%%%%%%%%%%%%%%%%%%%%%%%%%%%%%%%%%%%%%%
\subsection{Build-In Methods for Strings}
\begin{frame}[fragile, allowframebreaks]
\frametitle{Build-In Methods for Strings}
\begin{itemize}
 \item Most classes already come with a set of very useful methods
\lstset{language=shell}
\begin{lstlisting}[numbers=none, postbreak=, postbreak=]
$ irb --simple-prompt
>> String.instance_methods false
=> [:<=>, :==, :===, :eql?, :hash, :casecmp, :+, :*, :%, :[], :[]=, :insert, :length, :size, :bytesize, :empty?, :=~, :match, :succ, :succ!, :next, :next!, :upto, :index, :rindex, :replace, :clear, :chr, :getbyte, :setbyte, :byteslice, :to_i, :to_f, :to_s, :to_str, :inspect, :dump, :upcase, :downcase, :capitalize, :swapcase, :upcase!, :downcase!, :capitalize!, :swapcase!, :hex, :oct, :split, :lines, :bytes, :chars, :codepoints, :reverse, :reverse!, :concat, :<<, :prepend, :crypt, :intern, :to_sym, :ord, :include?, :start_with?, :end_with?, :scan, :ljust, :rjust, :center, :sub, :gsub, :chop, :chomp, :strip, :lstrip, :rstrip, :sub!, :gsub!, :chop!, :chomp!, :strip!, :lstrip!, :rstrip!, :tr, :tr_s, :delete, :squeeze, :count, :tr!, :tr_s!, :delete!, :squeeze!, :each_line, :each_byte, :each_char, :each_codepoint, :sum, :slice, :slice!, :partition, :rpartition, :encoding, :force_encoding, :b, :valid_encoding?, :ascii_only?, :unpack, :encode, :encode!, :to_r, :to_c]
>> exit
\end{lstlisting}

\item If you are not sure what one of these methods does you can use \texttt{ri} to look it up
\lstset{language=shell}
\begin{lstlisting}[numbers=none, escapechar=!]
$ ri -T String.size
\end{lstlisting}
\end{itemize}

\end{frame}
%%%%%%%%%%%%%%%%%%%%%%%%%%%%%%%%%%%%%%%%%%%%%%%%%%%%%%%%%%%%%%%%%%%%%%%%%%%%%%
%%%%%%%%%%%%%%%%%%%%%%%%%%%%%%%%%%%%%%%%%%%%%%%%%%%%%%%%%%%%%%%%%%%%%%%%%%%%%%
\subsection{Method Chaining}
\begin{frame}[fragile]
\frametitle{Method Chaining}
\begin{itemize}
 \item It is perfectly logical to chain different methods in Ruby

\lstset{language=Ruby, style=eclipse}
\begin{lstlisting}[escapechar=&]
# chaining.rb

a = String.new('A dog').upcase.reverse # Or String.new 'A dog'.upcase.reverse?
puts a
\end{lstlisting}

\lstset{language=shell}
\begin{lstlisting}[numbers=none, escapechar=!]
$ ruby chaining.rb
!\outputcommand{
GOD A
}!
\end{lstlisting}
\end{itemize}

\end{frame}
%%%%%%%%%%%%%%%%%%%%%%%%%%%%%%%%%%%%%%%%%%%%%%%%%%%%%%%%%%%%%%%%%%%%%%%%%%%%%%
%%%%%%%%%%%%%%%%%%%%%%%%%%%%%%%%%%%%%%%%%%%%%%%%%%%%%%%%%%%%%%%%%%%%%%%%%%%%%
\section{if-Condition}
\begin{frame}[fragile, allowframebreaks]
\frametitle{if-Condition}
\lstset{language=Ruby, style=eclipse}
\begin{lstlisting}[escapechar=&]
# if.rb

a=10

if a==10
  puts 'a is 10'
end

puts 'a is 10' if a==10
puts 'a is 10' unless a!=10

if a==10
  puts 'a is 10'
else
  puts 'a is not 10'
end

if a==10
  puts 'a is 10'
elsif a==20
  puts 'a is 20'
end
\end{lstlisting}

\lstset{language=shell}
\begin{lstlisting}[numbers=none, escapechar=!]
$ ruby getters_and_setters.rb
!\outputcommand{
a is 10\\
a is 10\\
a is 10\\
a is 10\\
a is 10\\
}!
\end{lstlisting}
\end{frame}
%%%%%%%%%%%%%%%%%%%%%%%%%%%%%%%%%%%%%%%%%%%%%%%%%%%%%%%%%%%%%%%%%%%%%%%%%%%%%%
%%%%%%%%%%%%%%%%%%%%%%%%%%%%%%%%%%%%%%%%%%%%%%%%%%%%%%%%%%%%%%%%%%%%%%%%%%%%%
\section{Loops}
\begin{frame}[fragile]
\frametitle{Loops}
\begin{itemize}
 \item There are different ways of implementing loops in Ruby. The iterator is used particularly often in the Rails environment
\end{itemize}

\end{frame}
%%%%%%%%%%%%%%%%%%%%%%%%%%%%%%%%%%%%%%%%%%%%%%%%%%%%%%%%%%%%%%%%%%%%%%%%%%%%%%
%%%%%%%%%%%%%%%%%%%%%%%%%%%%%%%%%%%%%%%%%%%%%%%%%%%%%%%%%%%%%%%%%%%%%%%%%%%%%
\subsection{While and until}
\begin{frame}[fragile]
\frametitle{While and until}
\lstset{language=Ruby, style=eclipse}
\begin{lstlisting}[escapechar=&]
# while.rb

i = 0
while i < 5 do
  puts i
  i = i+1
end
\end{lstlisting}

\lstset{language=shell}
\begin{lstlisting}[numbers=none, escapechar=!]
$ ruby while.rb
!\outputcommand{
0\\
1\\
2\\
3\\
4}!
\end{lstlisting}
\end{frame}
%%%%%%%%%%%%%%%%%%%%%%%%%%%%%%%%%%%%%%%%%%%%%%%%%%%%%%%%%%%%%%%%%%%%%%%%%%%%%%
%%%%%%%%%%%%%%%%%%%%%%%%%%%%%%%%%%%%%%%%%%%%%%%%%%%%%%%%%%%%%%%%%%%%%%%%%%%%%
\subsection{Block and Iterators}
\begin{frame}[fragile]
\frametitle{Block and Iterators}
\lstset{language=Ruby, style=eclipse}
\begin{lstlisting}[escapechar=&]
# times.rb

5.times { |i| puts i}

5.times do |i|
  puts i
end
\end{lstlisting}

\lstset{language=shell}
\begin{lstlisting}[numbers=none, escapechar=!]
$ ruby while.rb
!\outputcommand{
0\\
1\\
2\\
3\\
4\\
0\\
1\\
2\\
3\\
4}!
\end{lstlisting}
\end{frame}
%%%%%%%%%%%%%%%%%%%%%%%%%%%%%%%%%%%%%%%%%%%%%%%%%%%%%%%%%%%%%%%%%%%%%%%%%%%%%%
%%%%%%%%%%%%%%%%%%%%%%%%%%%%%%%%%%%%%%%%%%%%%%%%%%%%%%%%%%%%%%%%%%%%%%%%%%%%%
\section{Arrays and Hashes}
\begin{frame}[fragile, allowframebreaks]
\frametitle{Arrays and Hashes}
\begin{itemize}
 \item As in many programming languages, ``arrays'' and ``hashes'' are popular structures in Ruby for storing data
 \item An array is a list of objects 
\lstset{language=Ruby, style=eclipse}
\begin{lstlisting}[escapechar=&]
# arrays.rb

a = [3, 4 ,5]
a.each do |i|
  puts i
end

b = a.map{|i| i*i}
b.each{ |i| puts i}
\end{lstlisting}

\lstset{language=shell}
\begin{lstlisting}[numbers=none, escapechar=!]
$ ruby arrays.rb
!\outputcommand{
3\\
4\\
5\\
9\\
16\\
25}!
\end{lstlisting}
 \item A ``Hash'' is a list of key/value pairs
\lstset{language=Ruby, style=eclipse}
\begin{lstlisting}[escapechar=&]
# hashes.rb

price = {'eggs' => 0.99, 'butter' => 9.9}

puts price['eggs']
puts price['butter']

price.each do |key, value|
  puts "#{key}, #{value}"
end
\end{lstlisting}

\lstset{language=shell}
\begin{lstlisting}[numbers=none, escapechar=!]
$ ruby hashes.rb
!\outputcommand{
0.99\\
9.9\\
eggs, 0.99\\
butter, 9.9
}!
\end{lstlisting}

\item Notation for hashes (Ruby $>$ 2.0.0)
\lstset{language=Ruby, style=eclipse}
\begin{lstlisting}[escapechar=&]
# hashes.rb
price1 = {'egg' => 0.99, 'butter' => 9.9}
price2 = {:egg => 0.99, :butter => 9.9}
price3 = {egg: 0.99, butter: 9.9}

puts "price1 is#{(price1==price2)?' ': ' not '} equal to price2"
puts "price2 is#{(price2==price2)?' ': ' not '} equal to price3"
\end{lstlisting}

\lstset{language=shell}
\begin{lstlisting}[numbers=none, escapechar=!]
$ ruby hashes.rb
!\outputcommand{
price1 is not  equal to price2\\
price2 is equal to price3
}!
\end{lstlisting}
\end{itemize}

\end{frame}
%%%%%%%%%%%%%%%%%%%%%%%%%%%%%%%%%%%%%%%%%%%%%%%%%%%%%%%%%%%%%%%%%%%%%%%%%%%%%%
\section{Regular Expressions}
\begin{frame}[fragile, allowframebreaks]
\frametitle{Regular Expressions}
\begin{itemize}
 \item A regular expression is a pattern that can be matched against a string
 \item The most common way of creating a regular expression patter is to write it between forward slashes
 \item All characters except \texttt{., |, (, ), [, ], \{, \}, +, \, \^{}, \$, *, ? } match themselves
\end{itemize}
\pagebreak
 \begin{tabular}{ll}
/[Rr]uby/  & Match ``Ruby'' or ``ruby''\\
/rub[ye]/ & Match ``ruby'' or ``rube''\\
/[aeiou]/ & Match any one lowercase vowel\\
/[0-9]/ & Match any digit; same as /[0123456789]/\\
/[a-z]/   & Match any lowercase ASCII letter\\
/[A-Z]/   & Match any uppercase ASCII letter\\
/[a-zA-Z0-9] & Match any of the above\\
/[\^{}aeiou]/  & Match anything other than a lowercase vowel\\
/[\^{}0-9]/ & Match anything other than a digit\\
 \end{tabular} 
\end{frame}
%%%%%%%%%%%%%%%%%%%%%%%%%%%%%%%%%%%%%%%%%%%%%%%%%%%%%%%%%%%%%%%%%%%%%%%%%%%%%
\section{More? }
\begin{frame}[fragile]
\frametitle{More}
\begin{itemize}
 \item Symbols
 \item Mixing
 \item Exceptions, Catch and Throw
 \item Input and Output
 \item Metaprogramming
\end{itemize}
\end{frame}
%%%%%%%%%%%%%%%%%%%%%%%%%%%%%%%%%%%%%%%%%%%%%%%%%%%%%%%%%%%%%%%%%%%%%%%%%%%%%%

\end{document}
